\subsubsection{Descent of operations to the evaluation quotient}
\label{sec:ch-descenso-fibras}

The preceding quotient identifies histories publishing the same value, not
their complete registers. The following criterion characterizes exactly
when an operation on histories induces an operation on that quotient without
choosing a privileged history within each fiber; its formulation
continues the genealogical construction established in the preceding sections.

\begin{proposition}[Descent criterion through evaluation fibers]
\label{prop:ch-descenso-fibras}
Let \(q:\Omega\to J\) be the evaluation of a genealogical
chart onto its image. Let \(\star:D\to\Omega\) be an operation with
\(D\subseteq\Omega^2\) saturated by the fibers of \(q\times q\).
There exists a unique operation
\(\overline\star:(q\times q)(D)\to J\) such that
\[
 q(x\star y)=q(x)\,\overline\star\,q(y)
\]
if and only if, for pairs in the domain,
\[
 q(x)=q(x'),\quad q(y)=q(y')
 \quad\Longrightarrow\quad q(x\star y)=q(x'\star y').
\]
\end{proposition}

\begin{proof}
If the visible operation exists, its result depends only on the two
values, proving necessity. For sufficiency, define
\(s\,\overline\star\,t=q(x\star y)\), where \(q(x)=s\) and
\(q(y)=t\). Saturation preserves the domain when representatives are changed;
the stated condition preserves the result. Surjectivity of
\(q\times q\) onto the image of the domain proves uniqueness.
\end{proof}

A coordinate does not automatically replace a history as an argument of
all its operations: the criterion determines when that substitution is
legitimate. In particular, transporting an operation through a bijection
does not suffice to identify it with a native operation preceding the reading.
If a chart \(\Omega\) is compact and \(q\) is continuous and surjective
onto a Hausdorff space \(J\), evaluation also induces a
homeomorphism \(\Omega/{\sim_q}\to J\): it is a continuous bijection from a
compact space to a Hausdorff space. It is the \emph{quotient}, not necessarily the
space of histories, that is topologically identified with \(J\).
